% owned/owned.tex
% mainfile: ../perfbook.tex
% SPDX-License-Identifier: CC-BY-SA-3.0

\QuickQuizChapter{chp:Data Ownership}{数据所有权}{qqzowned}
%
\Epigraph{那是我的，我告诉你。
	  我自己的。
	  我的宝贝。
	  是的，我的宝贝。}
	 {Gollum，出自\emph{《魔戒：护戒使者》}，J.R.R.~Tolkien}

避免locking带来的同步开销，最简单的方法之一就是将数据分配给各个thread
（或者在内核的情况下，分配给各个CPU），
使得每一份数据只被一个thread访问和修改。
有趣的是，data ownership涵盖了三大并行设计技术中的每一种：
它在thread（或CPU，视情况而定）之间进行分区，
它将所有本地操作进行批处理，
而它对同步操作的消除则是弱化策略的逻辑极端。
因此，data ownership被大量使用就不足为奇了：
即使是新手也几乎会本能地使用它。
事实上，它的使用如此广泛，以至于本章不会引入任何新的示例，
而是回顾前面章节中的那些例子。

\QuickQuiz{
	在使用C或C++创建shared-memory并行程序
	（例如使用pthreads）时，哪种形式的data ownership
	极其难以避免？
}\QuickQuizAnswer{
	函数中auto变量的使用。
	默认情况下，这些变量是执行当前函数的thread的私有变量。
}\QuickQuizEnd

Data ownership有多种实现方式。
\Cref{sec:owned:Multiple Processes}介绍了data ownership的逻辑极端，
即每个thread拥有自己的私有地址空间。
\Cref{sec:owned:Partial Data Ownership and pthreads}探讨了
相反的极端，即数据是共享的，但不同的thread拥有对数据的不同访问权限。
\Cref{sec:owned:Function Shipping}描述了function shipping，
这是一种允许其他thread间接访问特定thread所拥有数据的方法。
\Cref{sec:owned:Designated Thread}描述了如何将指定thread
分配为特定功能及其相关数据的所有者。
\Cref{sec:owned:Privatization}讨论了通过将使用共享数据的算法
转换为使用data ownership来提升性能。
最后，\cref{sec:owned:Other Uses of Data Ownership}列举了
一些将data ownership作为一等公民的软件环境。

\section{多进程}
\label{sec:owned:Multiple Processes}
%
\epigraph{一个人的家就是他的城堡}{英格兰古法}

\Cref{sec:toolsoftrade:Scripting Languages}
介绍了以下示例：

\begin{VerbatimN}[samepage=true]
compute_it 1 > compute_it.1.out &
compute_it 2 > compute_it.2.out &
wait
cat compute_it.1.out
cat compute_it.2.out
\end{VerbatimN}

这个示例以并行方式运行两个\co{compute_it}程序实例，
作为不共享内存的独立进程。
因此，给定进程中的所有数据都归该进程所有，
所以上述示例中几乎所有的数据都是有所有权的。
这种方法几乎完全消除了同步开销。
由此带来的极致简洁性和最优性能的组合显然非常有吸引力。

\QuickQuizSeries{%
\QuickQuizB{
	在\cref{sec:owned:Multiple Processes}所示的示例中
	还残留着什么同步？
}\QuickQuizAnswerB{
	通过\co{sh}的\co{&}操作符创建thread，
	以及通过\co{sh}的\co{wait}命令汇合thread。

	当然，如果进程显式地共享内存，
	例如使用\co{shmget()}或\co{mmap()}系统调用，
	那么在访问或更新共享内存时很可能需要显式的同步。
	进程还可以使用以下任何进程间通信机制进行同步：
	\begin{enumerate}
	\item	System V信号量。
	\item	System V消息队列。
	\item	UNIX域套接字。
	\item	网络协议，包括TCP/IP、UDP以及大量其他协议。
	\item	文件锁定。
	\item	使用带有\co{O_CREAT}和\co{O_EXCL}标志的
		\co{open()}系统调用。
	\item	使用\co{rename()}系统调用。
	\end{enumerate}
	完整列出所有可能的同步机制留作读者的练习，
	在此提醒读者这将是一个极其冗长的列表。
	数量惊人的看似普通的系统调用都可以被用作同步机制。
}\QuickQuizEndB
%
\QuickQuizE{
	在\cref{sec:owned:Multiple Processes}所示的示例中
	存在共享数据吗？
}\QuickQuizAnswerE{
	这是一个哲学问题。

	希望答案是"否"的人可能会争辩说，
	按照定义，进程不共享内存。

	希望回答"是"的人可能会列出大量不需要共享内存的同步机制，
	指出内核会有一些共享状态，
	甚至可能会争辩说进程ID（PID）的分配构成了共享数据。

	这样的争论是很好的智力锻炼，
	也是展示聪明才智和在无助的同学或同事面前得分的好方法，
	但基本上只是一种避免做任何有用事情的方式。
}\QuickQuizEndE
}

这种模式同样可以用C语言编写，而不仅仅是\co{sh}，如
\cref{lst:toolsoftrade:Using the fork() Primitive,%
lst:toolsoftrade:Using the wait() Primitive}所示。

值得再次强调的是，这些简单的并行形式绝不是投机取巧或逃避责任，
而是使代码运行更快的简单而优雅的方法。
它速度快、可扩展性好、容易编程、容易维护，而且能完成任务。
此外，采用这种方法（在适用的情况下）可以让开发者有更多的时间
专注于其他事情，无论这些事情是对\co{compute_it}应用复杂的
单线程优化，还是对这种方法不适用的代码部分应用复杂的并行编程模式。
有什么理由不喜欢呢？

下一节讨论在shared-memory并行程序中使用data ownership。

\section{部分数据所有权与pthreads}
\label{sec:owned:Partial Data Ownership and pthreads}
%
\epigraph{与其关注你没有的，不如多关注你已拥有的。}
	 {Marcus Aurelius Antoninus}

并发计数（见\cref{chp:Counting}）大量使用了data ownership，
但增加了一个变化。
thread不允许修改其他thread拥有的数据，
但允许读取这些数据。
简而言之，shared memory的使用允许更细致的所有权和访问权限概念。

例如，考虑
\cref{lst:count:Per-Thread Statistical Counters}（位于
\cpageref{lst:count:Per-Thread Statistical Counters}）
中所示的per-thread统计计数器实现。
这里，\co{inc_count()}只更新对应thread的
\co{counter}实例，
而\co{read_count()}访问但不修改所有
thread的\co{counter}实例。

\QuickQuiz{
	是否有这样的部分data ownership场景：
	每个thread只读取自己的per-thread变量实例，
	但写入其他thread的实例？
}\QuickQuizAnswer{
	令人惊讶的是，确实有。
	一个例子是简单的消息传递系统，其中thread将消息发送到
	其他thread的邮箱中，每个thread负责在其发送的消息
	被处理后将其删除。
	实现这样一个算法留作读者的练习，
	识别具有类似所有权模式的其他算法也是如此。
}\QuickQuizEnd

部分data ownership在Linux内核中也很常见。
例如，某个CPU可能只有在禁用中断的情况下才被允许读取自己的
一组per-CPU变量，另一个CPU只有在持有相应的per-CPU lock时
才被允许读取第一个CPU的同一组per-CPU变量。
然后，该CPU在同时禁用中断并持有其per-CPU lock的情况下，
才被允许更新这组自己的per-CPU变量。
这种安排可以被看作一种reader-writer lock，它允许每个CPU
以非常低的开销访问自己的一组per-CPU变量。
这个主题有很多变体。

就其本身而言，纯粹的data ownership也既常见又有用，
例如在\cref{sec:SMPdesign:Resource Allocator Caches}
（从\cpageref{sec:SMPdesign:Resource Allocator Caches}开始）
中讨论的per-thread内存分配器cache。
在这个算法中，每个thread的cache完全是该thread私有的。

\section{功能转移}
\label{sec:owned:Function Shipping}
%
\epigraph{如果山不向穆罕默德走来，那么穆罕默德就必须走向山。}
	 {\emph{《随笔集》}，Francis Bacon}

上一节描述了一种弱形式的data ownership，其中thread伸手去获取
其他thread的数据。
这可以被理解为将数据带到需要它的函数那里。
另一种方法是将函数发送到数据所在之处。

这种方法的示例见
\cref{sec:count:Signal-Theft Limit Counter Design}
（从\cpageref{sec:count:Signal-Theft Limit Counter Design}开始），
特别是
\cref{lst:count:Signal-Theft Limit Counter Value-Migration Functions}
（位于
\cpageref{lst:count:Signal-Theft Limit Counter Value-Migration Functions}）
中的\co{flush_local_count_sig()}和
\co{flush_local_count()}函数。

\co{flush_local_count_sig()}函数是一个信号处理程序，
充当被转移的函数。
\co{flush_local_count()}中的\co{pthread_kill()}函数
发送信号——转移函数——然后等待被转移的函数执行完毕。
这个被转移的函数有一个常见的额外复杂性，
即需要与任何并发执行的\co{add_count()}
或\co{sub_count()}函数进行交互（见
\cref{lst:count:Signal-Theft Limit Counter Add Function}，
位于\cpageref{lst:count:Signal-Theft Limit Counter Add Function}，
以及\cref{lst:count:Signal-Theft Limit Counter Subtract Function}，
位于\cpageref{lst:count:Signal-Theft Limit Counter Subtract Function}）。

\QuickQuiz{
	除了POSIX信号之外，还有哪些机制可以用于function shipping？
}\QuickQuizAnswer{
	此类机制数量非常庞大，包括：
	\begin{enumerate}
	\item	System V消息队列。
	\item	Shared-memory双端队列（见
		\cref{sec:SMPdesign:Double-Ended Queue}）。
	\item	Shared-memory邮箱。
	\item	UNIX域套接字。
	\item	TCP/IP或UDP，可能辅以任何数量的高级协议，
		包括RPC、HTTP、XML、SOAP等等。
	\end{enumerate}
	编译完整列表留给足够专注的读者，
	在此提醒该列表将会非常冗长。
}\QuickQuizEnd

\section{指定线程}
\label{sec:owned:Designated Thread}
%
\epigraph{让一个人从事他最擅长的职业。}
	 {Cicero}

前面几节描述了允许每个thread保留自己的数据副本或数据部分的方法。
相比之下，本节描述了一种功能分解方法，
其中一个特殊的指定thread拥有完成其工作所需数据的权限。
\begin{fcvref}[ln:count:count_stat_eventual:whole:eventual]
\cref{sec:count:Eventually Consistent Implementation}中描述的
最终一致性计数器实现提供了一个例子。
该实现有一个运行\co{eventual()}函数的指定thread，
如\cref{lst:count:Array-Based Per-Thread Eventually Consistent Counters}
的\clnrefrange{b}{e}所示。
这个\co{eventual()} thread定期将per-thread计数值
拉入全局计数器中，这样对全局计数器的访问将如其名称所示，
最终收敛到实际值。
\end{fcvref}

\QuickQuiz{
	\begin{fcvref}[ln:count:count_stat_eventual:whole:eventual]
	但是\cref{lst:count:Array-Based Per-Thread Eventually Consistent Counters}
	的\clnrefrange{b}{e}所示的\co{eventual()}函数中的数据
	实际上都不归\co{eventual()} thread所有！
	这到底算什么data ownership？？？
	\end{fcvref}
}\QuickQuizAnswer{
	\begin{fcvref}[ln:count:count_stat_eventual:whole]
	关键短语是"拥有对数据的权限"。
	在这种情况下，所讨论的权限是访问在该列表
	\clnref{per_thr_cnt}定义的per-thread \co{counter}变量的权限。
	这种情况类似于
	\cref{sec:owned:Partial Data Ownership and pthreads}中描述的情况。

	然而，确实有归\co{eventual()} thread所有的数据，
	即在该列表\clnref{t,sum}定义的\co{t}和\co{sum}变量。

	有关指定thread的其他示例，请查看Linux内核中的内核thread，
	例如由\co{kthread_create()}和\co{kthread_run()}创建的那些。
	\end{fcvref}
}\QuickQuizEnd

\section{私有化}
\label{sec:owned:Privatization}
%
\epigraph{当然，一个人所夺取的和他真正拥有的之间是有区别的。}
	 {Pearl S.~Buck}

提高shared-memory并行程序性能和可扩展性的一种方法是
对其进行转换，将共享数据转换为由特定thread拥有的私有数据。

一个很好的例子见于\cref{sec:SMPdesign:Dining Philosophers Problem}
中一道Quick Quiz的答案，
该答案使用私有化为哲学家就餐问题提供了一个
比标准教科书解法具有更好性能和可扩展性的方案。
原始问题是五位哲学家围坐在桌旁，
每对相邻哲学家之间有一把叉子，最多允许两位哲学家同时进餐。

我们可以通过提供额外的五把叉子来轻松私有化这个问题，
这样每位哲学家都拥有自己的一对私有叉子。
这允许所有五位哲学家同时进餐，
同时也大大减少了某些类型疾病的传播。

在其他情况下，私有化会带来成本。
例如，考虑
\cref{lst:count:Simple Limit Counter Add; Subtract; and Read}（位于
\cpageref{lst:count:Simple Limit Counter Add; Subtract; and Read}）
中所示的简单限值计数器。
这是一个算法示例，其中thread可以读取彼此的数据，
但只允许更新自己的数据。
快速审查该算法可知，唯一的跨thread访问
在\co{read_count()}的求和循环中。
如果消除这个循环，我们就转向了更高效的纯data ownership，
但代价是\co{read_count()}的结果不太精确。

\QuickQuiz{
	是否可以在保持per-thread数据完全私有的同时
	获得更高的精度？
}\QuickQuizAnswer{
	是的。
	一种方法是让\co{read_count()}加上其自身per-thread变量的值。
	这保持了完全的所有权和性能，但精度只有轻微的提升，
	特别是在thread数量非常大的系统上。

	另一种方法是让\co{read_count()}使用function shipping，
	例如以per-thread信号的形式。
	这大大提高了精度，但\co{read_count()}的性能成本显著。

	然而，这两种方法都有一个优点，
	即消除了更新计数器这一常见情况下的cache抖动。
}\QuickQuizEnd

部分私有化也是可能的，有一些同步要求，
但比完全共享的情况要少。
\cref{sec:toolsoftrade:Avoiding Data Races}中探讨了
一些部分私有化的可能性。
\Cref{chp:Deferred Processing}将通过提供安全地将公共数据结构
变为私有的方法，为data ownership引入时间维度。

简而言之，私有化是并行程序员工具箱中的强大工具，
但仍然必须谨慎使用。
就像其他所有同步原语一样，它有可能在降低性能和可扩展性的同时
增加复杂性。

\section{数据所有权的其他用途}
\label{sec:owned:Other Uses of Data Ownership}
%
\epigraph{如果我们能创造接纳的能力，
	  属于我们的一切都会来到我们身边。}
	 {Rabindranath Tagore}

当数据可以被分区，使得跨thread访问或更新的需求很少或没有时，
data ownership效果最好。
幸运的是，这种情况相当常见，而且存在于各种各样的并行编程环境中。

Data ownership的例子包括：

\begin{enumerate}
\item	所有消息传递环境，如MPI~\cite{MPIForum2008}
	和BOINC~\cite{BOINC2008}。
\item	Map-reduce~\cite{MapReduce2008MIT}。
\item	客户端-服务器系统，包括RPC、Web服务，
	以及几乎所有带有后端数据库服务器的系统。
\item	Shared-nothing数据库系统。
\item	具有独立per-process地址空间的Fork-join系统。
\item	基于进程的并行，例如Erlang语言。
\item	私有变量，例如在thread环境中的C语言栈上auto变量。
\item	许多并行线性代数算法，尤其是那些非常适合GPGPU的算法。\footnote{
		但请注意，大量其他类别的应用程序也已被移植到
		GPGPU上~\cite{NormMatloff2017ParProcBook,AMD2020ROCm,NVidia2017GPGPU,NVidia2017GPGPU-university}。}
\item	为网络应用而改造的操作系统内核，其中每个连接
	（也称为\emph{流}~\cite{Shenker89,ZhangPhD,McKenney90}）
	被分配给特定的thread。
	这种方法的一个近期例子是IX操作系统~\cite{Belay:2016:IOS:3014162.2997641}。
	IX确实有一些共享数据结构，它们使用将在
	\cref{sec:defer:Read-Copy Update (RCU)}中描述的同步机制。
\end{enumerate}

Data ownership也许是现有最被低估的同步机制。
正确使用时，它提供了无与伦比的简洁性、性能和可扩展性。
也许正是它的简洁性使它没有得到应有的尊重。
希望对data ownership的微妙之处和威力有更深入的认识，
能够带来更高层次的尊重，更不用说带来更高的性能和可扩展性，
同时降低复杂性。

% populate with problems showing benefits of coupling data ownership with
% other approaches. For example, work-stealing schedulers. Perhaps also
% move memory allocation here, though its current location is quite good.

\QuickQuizAnswersChp{qqzowned}
